% --- Archon physics formalization source begin ---
% archon:physics
% archon:covers PhyXMiniProblems/problem_phyx_mini_0177.lean
% archon:source-report reports/phyx_mini/problem_phyx_mini_0177.source.json

\chapter{Physics problem phyx\_mini\_0177}
\label{ch:PhyXMiniProblems_problem_phyx_mini_0177}

\paragraph{Problem source.}
\#\# Physical scenario

Figure shows a standing wave on a \$2.0 \textbackslash{}, \textbackslash{}text\{m\}\$-long string that has been fixed at both ends and tightened until the wave speed is \$40 \textbackslash{}, \textbackslash{}text\{m/s\}\$.

\#\# Question

What is the frequency?

\#\# Answer choices

- A: A: 25 Hz

- B: B: 100 Hz

- C: C: 20 Hz

- D: D: 50 Hz

\#\# Figure caption (auxiliary; use the image as primary evidence)

The image shows a standing wave pattern, characterized by multiple nodes and antinodes along a straight line. The wave consists of symmetrical sine wave loops, and is depicted in blue. The edges on both sides feature gray gradient rectangles, likely representing boundaries or barriers. These gradients fade outwards, providing a visual boundary where the wave is contained. The straight line acts as the axis of the wave pattern, with alternating crests and troughs regularly spaced. The image conveys a clear demonstration of wave interference or resonance.

\#\# Dataset metadata

Resonance and Harmonics; Physical Model Grounding Reasoning, Multi-Formula Reasoning

\paragraph{Recorded answer/context.}
D: D: 50 Hz

\paragraph{Figure/image path.}
/root/proposal\_for\_physic/science-mango/phyx\_mini\_run/phyx\_data/test\_image/177.png

\paragraph{Formalization target.}
create a compiling Lean file with sorry bodies at `PhyXMiniProblems/problem\_phyx\_mini\_0177.lean`. The Lean declarations must preserve the physical quantities, dimensions or dimensional roles, figure labels, governing-law hypotheses, and final relation expressed by this problem.
Use Mathlib/Physlib names found through LeanExplore where available. If a physics API is missing, introduce faithful local abstractions rather than scalar placeholder aliases.

\begin{theorem}[Physics formalization target]
\label{thm:physics:phyx_mini_0177:target}
\lean{PhyXMiniProblems.ProblemPhyXMini0177.problem_phyx_mini_0177}
\uses{def:physics:phyx-mini-0177:phyxminiproblems-problemphyxmini0177-frequencyinhertz, def:physics:phyx-mini-0177:phyxminiproblems-problemphyxmini0177-fixedstringstandingwavesetup, def:physics:phyx-mini-0177:phyxminiproblems-problemphyxmini0177-matchesfixedstringscenario, def:physics:phyx-mini-0177:phyxminiproblems-problemphyxmini0177-matchesproblemreadouts, def:physics:phyx-mini-0177:phyxminiproblems-problemphyxmini0177-matchessuppliedstandingwavefigure, def:physics:phyx-mini-0177:phyxminiproblems-problemphyxmini0177-hasphysicalstandingwaveparameters, def:physics:phyx-mini-0177:phyxminiproblems-problemphyxmini0177-satisfiesfixedendstandingwavemodelaw, def:physics:phyx-mini-0177:phyxminiproblems-problemphyxmini0177-satisfiesstringwavespeedlaw, def:physics:phyx-mini-0177:phyxminiproblems-problemphyxmini0177-recordeddatasetanswer, def:physics:phyx-mini-0177:phyxminiproblems-problemphyxmini0177-isexactfrequencychoice}
The assigned autoformalize agent should translate this physics problem into Lean declarations in the covered file, with theorem and lemma proof bodies written as `by sorry`.
\end{theorem}
\begin{proof}
This is an autoformalization task, not a proof task. Produce statements that can later be proved by the physics prover without weakening the physical model.
\end{proof}

% --- Archon declaration topology begin ---
\section{Declaration topology}
\begin{definition}[Length Quantity]
  \label{def:physics:phyx-mini-0177:phyxminiproblems-problemphyxmini0177-lengthquantity}
  \lean{PhyXMiniProblems.ProblemPhyXMini0177.LengthQuantity}
  A nonnegative physical length, independent of the unit used to read it
\end{definition}
\begin{proof}
  This is the declared model definition of Length Quantity.
\end{proof}
\begin{definition}[Frequency Quantity]
  \label{def:physics:phyx-mini-0177:phyxminiproblems-problemphyxmini0177-frequencyquantity}
  \lean{PhyXMiniProblems.ProblemPhyXMini0177.FrequencyQuantity}
  A nonnegative physical frequency, carrying inverse-time dimension
\end{definition}
\begin{proof}
  This is the declared model definition of Frequency Quantity.
\end{proof}
\begin{definition}[Speed Quantity]
  \label{def:physics:phyx-mini-0177:phyxminiproblems-problemphyxmini0177-speedquantity}
  \lean{PhyXMiniProblems.ProblemPhyXMini0177.SpeedQuantity}
  A nonnegative physical wave-propagation speed
\end{definition}
\begin{proof}
  This is the declared model definition of Speed Quantity.
\end{proof}
\begin{definition}[length Readout]
  \label{def:physics:phyx-mini-0177:phyxminiproblems-problemphyxmini0177-lengthreadout}
  \lean{PhyXMiniProblems.ProblemPhyXMini0177.lengthReadout}
  \uses{def:physics:phyx-mini-0177:phyxminiproblems-problemphyxmini0177-lengthquantity}
  Read a physical length in a selected length unit
\end{definition}
\begin{proof}
  This is the declared model definition of length Readout.
\end{proof}
\begin{definition}[frequency Readout]
  \label{def:physics:phyx-mini-0177:phyxminiproblems-problemphyxmini0177-frequencyreadout}
  \lean{PhyXMiniProblems.ProblemPhyXMini0177.frequencyReadout}
  \uses{def:physics:phyx-mini-0177:phyxminiproblems-problemphyxmini0177-frequencyquantity}
  Read a physical frequency in inverse units of a selected time unit
\end{definition}
\begin{proof}
  This is the declared model definition of frequency Readout.
\end{proof}
\begin{definition}[speed Readout]
  \label{def:physics:phyx-mini-0177:phyxminiproblems-problemphyxmini0177-speedreadout}
  \lean{PhyXMiniProblems.ProblemPhyXMini0177.speedReadout}
  \uses{def:physics:phyx-mini-0177:phyxminiproblems-problemphyxmini0177-speedquantity}
  Read a physical speed in selected length and time units
\end{definition}
\begin{proof}
  This is the declared model definition of speed Readout.
\end{proof}
\begin{definition}[length In Meters]
  \label{def:physics:phyx-mini-0177:phyxminiproblems-problemphyxmini0177-lengthinmeters}
  \lean{PhyXMiniProblems.ProblemPhyXMini0177.lengthInMeters}
  \uses{def:physics:phyx-mini-0177:phyxminiproblems-problemphyxmini0177-lengthquantity, def:physics:phyx-mini-0177:phyxminiproblems-problemphyxmini0177-lengthreadout}
  Meter readout of a physical length
\end{definition}
\begin{proof}
  This is the declared model definition of length In Meters.
\end{proof}
\begin{definition}[frequency In Hertz]
  \label{def:physics:phyx-mini-0177:phyxminiproblems-problemphyxmini0177-frequencyinhertz}
  \lean{PhyXMiniProblems.ProblemPhyXMini0177.frequencyInHertz}
  \uses{def:physics:phyx-mini-0177:phyxminiproblems-problemphyxmini0177-frequencyquantity, def:physics:phyx-mini-0177:phyxminiproblems-problemphyxmini0177-frequencyreadout}
  Hertz readout of a physical frequency
\end{definition}
\begin{proof}
  This is the declared model definition of frequency In Hertz.
\end{proof}
\begin{definition}[speed In Meters Per Second]
  \label{def:physics:phyx-mini-0177:phyxminiproblems-problemphyxmini0177-speedinmeterspersecond}
  \lean{PhyXMiniProblems.ProblemPhyXMini0177.speedInMetersPerSecond}
  \uses{def:physics:phyx-mini-0177:phyxminiproblems-problemphyxmini0177-speedquantity, def:physics:phyx-mini-0177:phyxminiproblems-problemphyxmini0177-speedreadout}
  Meter-per-second readout of a physical speed
\end{definition}
\begin{proof}
  This is the declared model definition of speed In Meters Per Second.
\end{proof}
\begin{definition}[String Endpoint]
  \label{def:physics:phyx-mini-0177:phyxminiproblems-problemphyxmini0177-stringendpoint}
  \lean{PhyXMiniProblems.ProblemPhyXMini0177.StringEndpoint}
  The two endpoints visible against the gray supports in the figure
\end{definition}
\begin{proof}
  This is the declared model definition of String Endpoint.
\end{proof}
\begin{definition}[Endpoint Boundary Condition]
  \label{def:physics:phyx-mini-0177:phyxminiproblems-problemphyxmini0177-endpointboundarycondition}
  \lean{PhyXMiniProblems.ProblemPhyXMini0177.EndpointBoundaryCondition}
  Boundary condition specified at either endpoint of the string
\end{definition}
\begin{proof}
  This is the declared model definition of Endpoint Boundary Condition.
\end{proof}
\begin{definition}[String Mechanical State]
  \label{def:physics:phyx-mini-0177:phyxminiproblems-problemphyxmini0177-stringmechanicalstate}
  \lean{PhyXMiniProblems.ProblemPhyXMini0177.StringMechanicalState}
  Mechanical state of the string described in the problem
\end{definition}
\begin{proof}
  This is the declared model definition of String Mechanical State.
\end{proof}
\begin{definition}[Transverse Wave Kind]
  \label{def:physics:phyx-mini-0177:phyxminiproblems-problemphyxmini0177-transversewavekind}
  \lean{PhyXMiniProblems.ProblemPhyXMini0177.TransverseWaveKind}
  Kind of transverse wave depicted by the blue curves
\end{definition}
\begin{proof}
  This is the declared model definition of Transverse Wave Kind.
\end{proof}
\begin{definition}[Standing Wave Figure]
  \label{def:physics:phyx-mini-0177:phyxminiproblems-problemphyxmini0177-standingwavefigure}
  \lean{PhyXMiniProblems.ProblemPhyXMini0177.StandingWaveFigure}
  \uses{def:physics:phyx-mini-0177:phyxminiproblems-problemphyxmini0177-stringendpoint}
  Discrete information read from the supplied standing-wave diagram. Nodes include both endpoints. An antinode count is equivalently the number of loops of a fixed-fixed normal mode
\end{definition}
\begin{proof}
  This is the declared model definition of Standing Wave Figure.
\end{proof}
\begin{definition}[Fixed String Standing Wave Setup]
  \label{def:physics:phyx-mini-0177:phyxminiproblems-problemphyxmini0177-fixedstringstandingwavesetup}
  \lean{PhyXMiniProblems.ProblemPhyXMini0177.FixedStringStandingWaveSetup}
  \uses{def:physics:phyx-mini-0177:phyxminiproblems-problemphyxmini0177-lengthquantity, def:physics:phyx-mini-0177:phyxminiproblems-problemphyxmini0177-frequencyquantity, def:physics:phyx-mini-0177:phyxminiproblems-problemphyxmini0177-speedquantity, def:physics:phyx-mini-0177:phyxminiproblems-problemphyxmini0177-stringendpoint, def:physics:phyx-mini-0177:phyxminiproblems-problemphyxmini0177-endpointboundarycondition, def:physics:phyx-mini-0177:phyxminiproblems-problemphyxmini0177-stringmechanicalstate, def:physics:phyx-mini-0177:phyxminiproblems-problemphyxmini0177-transversewavekind, def:physics:phyx-mini-0177:phyxminiproblems-problemphyxmini0177-standingwavefigure}
  The physical string and the particular normal mode shown in the problem. The wavelength and oscillation frequency are unknown dimensionful quantities; they are constrained only by the governing-law hypotheses below
\end{definition}
\begin{proof}
  This is the declared model definition of Fixed String Standing Wave Setup.
\end{proof}
\begin{definition}[Matches Fixed String Scenario]
  \label{def:physics:phyx-mini-0177:phyxminiproblems-problemphyxmini0177-matchesfixedstringscenario}
  \lean{PhyXMiniProblems.ProblemPhyXMini0177.MatchesFixedStringScenario}
  \uses{def:physics:phyx-mini-0177:phyxminiproblems-problemphyxmini0177-fixedstringstandingwavesetup}
  Qualitative information from the prose: the string is tightened, the wave is standing, and both ends are fixed. This predicate contains no wavelength or frequency value
\end{definition}
\begin{proof}
  This is the declared model definition of Matches Fixed String Scenario.
\end{proof}
\begin{definition}[Matches Problem Readouts]
  \label{def:physics:phyx-mini-0177:phyxminiproblems-problemphyxmini0177-matchesproblemreadouts}
  \lean{PhyXMiniProblems.ProblemPhyXMini0177.MatchesProblemReadouts}
  \uses{def:physics:phyx-mini-0177:phyxminiproblems-problemphyxmini0177-lengthinmeters, def:physics:phyx-mini-0177:phyxminiproblems-problemphyxmini0177-speedinmeterspersecond, def:physics:phyx-mini-0177:phyxminiproblems-problemphyxmini0177-fixedstringstandingwavesetup}
  Numerical quantities stated in the prose. They are read in SI units but remain stored as unit-independent physical quantities
\end{definition}
\begin{proof}
  This is the declared model definition of Matches Problem Readouts.
\end{proof}
\begin{definition}[Matches Supplied Standing Wave Figure]
  \label{def:physics:phyx-mini-0177:phyxminiproblems-problemphyxmini0177-matchessuppliedstandingwavefigure}
  \lean{PhyXMiniProblems.ProblemPhyXMini0177.MatchesSuppliedStandingWaveFigure}
  \uses{def:physics:phyx-mini-0177:phyxminiproblems-problemphyxmini0177-fixedstringstandingwavesetup}
  Primary-image evidence: both supported endpoints and the equilibrium axis are shown, and the blue pattern has five loops and six nodes including the ends. No frequency or wavelength readout is present in the bitmap
\end{definition}
\begin{proof}
  This is the declared model definition of Matches Supplied Standing Wave Figure.
\end{proof}
\begin{definition}[Has Physical Standing Wave Parameters]
  \label{def:physics:phyx-mini-0177:phyxminiproblems-problemphyxmini0177-hasphysicalstandingwaveparameters}
  \lean{PhyXMiniProblems.ProblemPhyXMini0177.HasPhysicalStandingWaveParameters}
  \uses{def:physics:phyx-mini-0177:phyxminiproblems-problemphyxmini0177-lengthinmeters, def:physics:phyx-mini-0177:phyxminiproblems-problemphyxmini0177-frequencyinhertz, def:physics:phyx-mini-0177:phyxminiproblems-problemphyxmini0177-speedinmeterspersecond, def:physics:phyx-mini-0177:phyxminiproblems-problemphyxmini0177-fixedstringstandingwavesetup}
  Positivity and nondegeneracy conditions for the physical normal mode
\end{definition}
\begin{proof}
  This is the declared model definition of Has Physical Standing Wave Parameters.
\end{proof}
\begin{definition}[Satisfies Fixed End Standing Wave Mode Law]
  \label{def:physics:phyx-mini-0177:phyxminiproblems-problemphyxmini0177-satisfiesfixedendstandingwavemodelaw}
  \lean{PhyXMiniProblems.ProblemPhyXMini0177.SatisfiesFixedEndStandingWaveModeLaw}
  \uses{def:physics:phyx-mini-0177:phyxminiproblems-problemphyxmini0177-lengthreadout, def:physics:phyx-mini-0177:phyxminiproblems-problemphyxmini0177-fixedstringstandingwavesetup}
  For a string fixed at both ends, normal mode n contains n antinodes and n + 1 nodes, while its length contains n half-wavelengths. The last equation is stated in every length unit and does not single out the pictured mode or its numerical wavelength
\end{definition}
\begin{proof}
  This is the declared model definition of Satisfies Fixed End Standing Wave Mode Law.
\end{proof}
\begin{definition}[Satisfies String Wave Speed Law]
  \label{def:physics:phyx-mini-0177:phyxminiproblems-problemphyxmini0177-satisfiesstringwavespeedlaw}
  \lean{PhyXMiniProblems.ProblemPhyXMini0177.SatisfiesStringWaveSpeedLaw}
  \uses{def:physics:phyx-mini-0177:phyxminiproblems-problemphyxmini0177-lengthreadout, def:physics:phyx-mini-0177:phyxminiproblems-problemphyxmini0177-frequencyreadout, def:physics:phyx-mini-0177:phyxminiproblems-problemphyxmini0177-speedreadout, def:physics:phyx-mini-0177:phyxminiproblems-problemphyxmini0177-fixedstringstandingwavesetup}
  The nondispersive wave relation v = lambda f, expressed in every compatible choice of length and time units. It is a general governing law, not the requested numerical frequency
\end{definition}
\begin{proof}
  This is the declared model definition of Satisfies String Wave Speed Law.
\end{proof}
\begin{definition}[Answer Choice]
  \label{def:physics:phyx-mini-0177:phyxminiproblems-problemphyxmini0177-answerchoice}
  \lean{PhyXMiniProblems.ProblemPhyXMini0177.AnswerChoice}
  Labels of the four frequency choices printed with the problem
\end{definition}
\begin{proof}
  This is the declared model definition of Answer Choice.
\end{proof}
\begin{definition}[displayed Frequency In Hertz]
  \label{def:physics:phyx-mini-0177:phyxminiproblems-problemphyxmini0177-displayedfrequencyinhertz}
  \lean{PhyXMiniProblems.ProblemPhyXMini0177.displayedFrequencyInHertz}
  \uses{def:physics:phyx-mini-0177:phyxminiproblems-problemphyxmini0177-answerchoice}
  Hertz value displayed beside each answer label
\end{definition}
\begin{proof}
  This is the declared model definition of displayed Frequency In Hertz.
\end{proof}
\begin{definition}[recorded Dataset Answer]
  \label{def:physics:phyx-mini-0177:phyxminiproblems-problemphyxmini0177-recordeddatasetanswer}
  \lean{PhyXMiniProblems.ProblemPhyXMini0177.recordedDatasetAnswer}
  \uses{def:physics:phyx-mini-0177:phyxminiproblems-problemphyxmini0177-answerchoice}
  Answer label recorded in the source dataset
\end{definition}
\begin{proof}
  This is the declared model definition of recorded Dataset Answer.
\end{proof}
\begin{definition}[Is Exact Frequency Choice]
  \label{def:physics:phyx-mini-0177:phyxminiproblems-problemphyxmini0177-isexactfrequencychoice}
  \lean{PhyXMiniProblems.ProblemPhyXMini0177.IsExactFrequencyChoice}
  \uses{def:physics:phyx-mini-0177:phyxminiproblems-problemphyxmini0177-frequencyinhertz, def:physics:phyx-mini-0177:phyxminiproblems-problemphyxmini0177-fixedstringstandingwavesetup, def:physics:phyx-mini-0177:phyxminiproblems-problemphyxmini0177-answerchoice, def:physics:phyx-mini-0177:phyxminiproblems-problemphyxmini0177-displayedfrequencyinhertz}
  A choice matches the physical frequency exactly when its Hertz readout does
\end{definition}
\begin{proof}
  This is the declared model definition of Is Exact Frequency Choice.
\end{proof}
\begin{lemma}[harmonic Index eq five]
  \label{lem:physics:phyx-mini-0177:phyxminiproblems-problemphyxmini0177-harmonicindex-eq-five}
  \lean{PhyXMiniProblems.ProblemPhyXMini0177.harmonicIndex_eq_five}
  \uses{def:physics:phyx-mini-0177:phyxminiproblems-problemphyxmini0177-fixedstringstandingwavesetup, def:physics:phyx-mini-0177:phyxminiproblems-problemphyxmini0177-matchessuppliedstandingwavefigure, def:physics:phyx-mini-0177:phyxminiproblems-problemphyxmini0177-satisfiesfixedendstandingwavemodelaw}
  The five loops in the supplied image identify the fifth fixed-fixed normal mode. This is a consequence of figure evidence and the general mode-count law, rather than a field initialized to the answer
\end{lemma}
\begin{proof}
  The conclusion follows from the governing relations and figure readouts named in the statement of harmonic Index eq five.
\end{proof}
\begin{lemma}[wavelength in meters eq four fifths]
  \label{lem:physics:phyx-mini-0177:phyxminiproblems-problemphyxmini0177-wavelength-in-meters-eq-four-fifths}
  \lean{PhyXMiniProblems.ProblemPhyXMini0177.wavelength_in_meters_eq_four_fifths}
  \uses{def:physics:phyx-mini-0177:phyxminiproblems-problemphyxmini0177-lengthinmeters, def:physics:phyx-mini-0177:phyxminiproblems-problemphyxmini0177-fixedstringstandingwavesetup, def:physics:phyx-mini-0177:phyxminiproblems-problemphyxmini0177-matchesfixedstringscenario, def:physics:phyx-mini-0177:phyxminiproblems-problemphyxmini0177-matchesproblemreadouts, def:physics:phyx-mini-0177:phyxminiproblems-problemphyxmini0177-matchessuppliedstandingwavefigure, def:physics:phyx-mini-0177:phyxminiproblems-problemphyxmini0177-hasphysicalstandingwaveparameters, def:physics:phyx-mini-0177:phyxminiproblems-problemphyxmini0177-satisfiesfixedendstandingwavemodelaw}
  For a two-meter string in its fifth mode, the fixed-end wavelength relation gives lambda = 4/5 m
\end{lemma}
\begin{proof}
  The conclusion follows from the governing relations and figure readouts named in the statement of wavelength in meters eq four fifths.
\end{proof}
\begin{lemma}[frequency in hertz eq fifty]
  \label{lem:physics:phyx-mini-0177:phyxminiproblems-problemphyxmini0177-frequency-in-hertz-eq-fifty}
  \lean{PhyXMiniProblems.ProblemPhyXMini0177.frequency_in_hertz_eq_fifty}
  \uses{def:physics:phyx-mini-0177:phyxminiproblems-problemphyxmini0177-frequencyinhertz, def:physics:phyx-mini-0177:phyxminiproblems-problemphyxmini0177-fixedstringstandingwavesetup, def:physics:phyx-mini-0177:phyxminiproblems-problemphyxmini0177-matchesfixedstringscenario, def:physics:phyx-mini-0177:phyxminiproblems-problemphyxmini0177-matchesproblemreadouts, def:physics:phyx-mini-0177:phyxminiproblems-problemphyxmini0177-matchessuppliedstandingwavefigure, def:physics:phyx-mini-0177:phyxminiproblems-problemphyxmini0177-hasphysicalstandingwaveparameters, def:physics:phyx-mini-0177:phyxminiproblems-problemphyxmini0177-satisfiesfixedendstandingwavemodelaw, def:physics:phyx-mini-0177:phyxminiproblems-problemphyxmini0177-satisfiesstringwavespeedlaw}
  Combining the derived 4/5 m wavelength with the stated 40 m/s speed in v = lambda f yields the exact frequency 50 Hz
\end{lemma}
\begin{proof}
  The conclusion follows from the governing relations and figure readouts named in the statement of frequency in hertz eq fifty.
\end{proof}
% --- Archon declaration topology end ---
% --- Archon physics formalization source end ---
